\documentclass[11pt]{article}
\usepackage[a4paper,margin=1in]{geometry}
\usepackage[T1]{fontenc}
\usepackage{lmodern}
\usepackage{amsmath,amssymb,amsthm,mathtools,booktabs}
\usepackage[numbers,sort&compress]{natbib}
\usepackage[colorlinks=true,allcolors=blue]{hyperref}
\hypersetup{pdfauthor={},pdftitle={Integer STU flux families: positive masses and tadpole-controlled coupling}}
\newtheorem{theorem}{Theorem}
\newtheorem{lemma}{Lemma}
\newtheorem{proposition}{Proposition}
\newcommand{\ii}{\mathrm{i}}
\newcommand{\R}{\mathbb R}
\newcommand{\Z}{\mathbb Z}
\newcommand{\Imn}{\operatorname{Im}}
\newcommand{\diag}{\operatorname{diag}}
\author{}
\date{29 October 2026}
\title{Integer STU flux families:\\positive masses and tadpole-controlled coupling}
\begin{document}
\maketitle
\begin{abstract}
An exact even integral STU flux lattice admits a positive-divisor family of
no-scale complex-structure/axiodilaton critical points. With $A_IB_I=k$, its
charge is $Q=8k$, its full real F-term Jacobian has determinant $-4096k^4$,
and its eight canonical masses squared are $k$ times a fixed positive
spectrum. An integral modular image gives $S=\ii k^2$, $U_i=\ii k$, hence
$g_s=k^{-2}$. At inherited quotient tadpole capacity $L=28$, the construction
reaches $g_s=1/9$ with $Q=24$ and contains nine distinct reduced moduli point
types. We also prove exact pointwise charge minima at those points using
primitive complementary-monomial pairs and a scaling obstruction. The
results are theorems of the stated period/lattice model. Its global
orientifold interpretation is inherited, and the nonzero superpotential,
unstabilized K\"ahler sector and common volume factor remain explicit.
\end{abstract}

\section{The exact model and its interpretation}
Flux stabilization and its source-charge budget are related but distinct
parts of the no-scale compactification problem
\citep{GKP,DenefEtAl,BenaEtAl}. Here we solve a specified low-dimensional
integral period problem; no general tadpole conjecture is proved or refuted.
Write $z=(S,U_1,U_2,U_3)$, with $y_A=\Imn z_A>0$, and take
\begin{align}
\Pi&=(1,U_1,U_2,U_3,U_1U_2U_3,-U_2U_3,-U_1U_3,-U_1U_2)^T,\label{period}\\
P(n)&=n_0+\sum_i n_iU_i-n^0U_1U_2U_3
       +n^1U_2U_3+n^2U_1U_3+n^3U_1U_2,\\
W&=P(f)-SP(h),\qquad K=-\sum_{A=0}^3\log(2y_A).
\end{align}
The flux ordering is $(f^I,f_I;h^I,h_I)$, $I=0,1,2,3$, with
$(f,h)\in(2\Z)^{16}$. Fix the quotient normalization
\begin{equation}
Q=\frac12\sum_I(f_Ih^I-f^Ih_I),\qquad N_{\rm cover}=2Q.
\label{charge}
\end{equation}
All results below use this precise sign and factor of two.

The geometric motivation is the exact untwisted three-complex-structure
period sector of the resolved toroidal model associated with
\citet{DenefEtAl}. The smooth orientifold realization, the interpretation of
this even sublattice as a conservative non-exotic O3-minus sector, and the
source budgets $L=28$ and $L=4$ are inherited from the preceding source
calculations. This note does not independently establish those global
geometric or admissibility assertions. Its algebraic claims hold for
\eqref{period}--\eqref{charge} whether or not one inserts either capacity.

An F-term zero means $D_AW=0$ for these four complex variables. We use the
eight-by-eight real Jacobian $J_F$ of $(\Re D_AW,\Im D_AW)$, in interleaved
real coordinates $(\Re z_0,\Im z_0,\ldots,\Re z_3,\Im z_3)$. This is the
rank relevant to the no-scale potential; a holomorphic Hessian alone is
insufficient when $W\ne0$.

\section{An integral divisor family}
\begin{theorem}[Construction and masses]\label{family}
Let $k>0$ and $A_I,B_I$ be positive integers satisfying $A_IB_I=k$ for every
$I$. Set $f^I=h_I=0$, $f_I=-2A_I$, $h^I=-2B_I$. Then
\begin{equation}
 U_i=\ii\frac{A_0}{A_i},\qquad
 S=\ii\frac{A_1A_2A_3}{A_0k}
 \label{divisorpoint}
\end{equation}
is an F-term zero with $Q=8k$ and $\det J_F=-4096k^4$. With the common
K\"ahler-sector factor $e^{K_K}$ set to one, its canonical masses squared are
\begin{equation}
k\{28-12\sqrt5,\ 28+12\sqrt5,\
 (12-4\sqrt5)^{[3]},\ (12+4\sqrt5)^{[3]}\}.
\label{spectrum}
\end{equation}
The superscript $[3]$ denotes multiplicity three. All masses squared are
strictly positive.
\end{theorem}

\begin{proof}
The unit seed has $f_I=h^I=-2$, other entries zero, and
\begin{equation}
W_0=-2(1+U_1+U_2+U_3)
 -2S(U_1U_2U_3-U_1U_2-U_1U_3-U_2U_3).
\end{equation}
At $z_A=\ii$, $W_0=-4-12\ii$ and every holomorphic first derivative equals
$-6+2\ii$. Since $K_A=\ii/2$ there, all four F-terms vanish. Direct real
differentiation gives the seed Jacobian in Appendix~\ref{seed}, with
determinant $-4096$.

For $c=A_0$ and the positive imaginary parts $y_A$ in
\eqref{divisorpoint}, coefficient matching gives
\begin{equation}
W(z)=cW_0(z_0/y_0,z_1/y_1,z_2/y_2,z_3/y_3),\qquad
\prod_Ay_A=A_0^2/k.
\label{covariance}
\end{equation}
The logarithmic K\"ahler derivatives imply
$D_AW=(c/y_A)(D_AW_0)(z/y)$. Thus the rescaled seed is a zero. The divisor
conditions ensure that all flux entries were integers before performing
this construction; no continuous flux is rounded. Equation~\eqref{charge}
gives $Q=\frac12\sum_I4A_IB_I=8k$.

Rows and columns of the real Jacobian scale separately, so
\begin{equation}
\det J_F=\frac{c^8}{(\prod_Ay_A)^4}\det J_{F,0}=-4096k^4.
\end{equation}
At an F-term zero the real potential Hessian and kinetic metric are
\begin{align}
H_{\rm real}&=2e^KJ_F^T\diag(4y_0^2,4y_0^2,\ldots,4y_3^2,4y_3^2)J_F,
\label{hessian}\\
G_{\rm real}&=\diag(1/(2y_0^2),1/(2y_0^2),\ldots,1/(2y_3^2),1/(2y_3^2)).
\end{align}
Here $\mathcal L_{\rm kin}=-\frac12G_{ab}\partial\phi^a\partial\phi^b$.
Canonical normalization of \eqref{hessian} and \eqref{covariance} gives
\begin{equation}
H_{\rm can}=\frac{c^2}{\prod_Ay_A}H_0=kH_0
\end{equation}
in aligned frames. At the unit seed,
$H_0=J_{F,0}^TJ_{F,0}$ and
\begin{equation}
\det(tI-H_0)=(t^2-56t+64)(t^2-24t+64)^3.
\end{equation}
Its roots give \eqref{spectrum}. Positivity follows, for example, from
$28^2>144\cdot5$ and $12^2>16\cdot5$.
\end{proof}

The simple subfamily $A_I=m$, $B_I=n$, $k=mn$ gives
$S=\ii m/n$, $U_i=\ii$, $Q=8mn$. A common integer multiple of $(m,n)$
changes the charge while leaving the moduli point fixed; it is not a new
moduli point. More generally \eqref{covariance} is a map between different
integral models after the divisor condition is imposed. It is not a
continuous family of fixed-flux isolated vacua or an assumed integral
duality.

\section{A weak-coupling branch and modular point types}
Choosing $A_0=1$, $A_i=k$ in Theorem~\ref{family}, then inverting all three
torus moduli, gives the explicit even flux
\begin{equation}
\begin{split}
 f^I&=(2,2k,2k,2k),\quad f_I=h^I=(0,0,0,0),\\
 h_I&=(-2k,-2,-2,-2),\qquad S=\ii k^2,\quad U_i=\ii k.
\end{split}
\label{weak}
\end{equation}
It can also be checked directly in
\begin{equation}
W=-2U_1U_2U_3+2k\sum_{i<j}U_iU_j+2S(k+U_1+U_2+U_3).
\end{equation}
At \eqref{weak},
\begin{equation}
W=4k^3(-3+\ii),\qquad Q=8k,\qquad g_s=k^{-2}.
\end{equation}
The determinant and spectrum remain those in Theorem~\ref{family}.
Capacity imposes $k\le\lfloor L/8\rfloor$ on this branch. Thus $L=28$
permits $k=1,2,3$, with respective couplings $1,1/4,1/9$ and residual
quotient D3 charges $20,12,4$. At $k=3$ the flux is
$(2,6,6,6;0^4;0^4;-6,-2,-2,-2)$, $N_{\rm cover}=48$, and $Q=24$.
Comparing 48 directly to capacity 28 would use inconsistent normalizations.

Along this branch $g_s=(8/Q)^2$. The weaker coupling at finite allowed
charge is a useful construction, but a parametrically weak-coupling limit
requires increasing capacity. Within the whole positive divisor family,
$k^{-2}\le\Imn S\le k^2$, so modular reduction cannot improve on
$g_s=1/k^2$ in that branch. Optimality over other even integral fluxes is
not established.
For $L\ge8$, the attainable branch minimum is precisely
$g_s=\lfloor L/8\rfloor^{-2}$, realized by its largest allowed integer $k$.

For clarity, separate new moduli points from their modular images. With
$\gamma=\left(\begin{smallmatrix}a&b\\c&d\end{smallmatrix}\right)\in
SL(2,\Z)$,
\begin{equation}
S'=\frac{aS+b}{cS+d},\quad f'=af+bh,\quad h'=cf+dh,\quad
W'=\frac{W}{cS+d}.
\end{equation}
These images preserve $Q$, even integrality, F-term zeros and masses, and
are not counted again. Torus inversion is the actual integer polynomial
map $W'(U_j')=U_j'W(-1/U_j')$; torus permutations relabel the $U_i$.

For $k=1,2,3$ there are $1+16+16=33$ divisor assignments. Reducing each
pure imaginary modulus by inversion and sorting the three torus factors
gives the nine point types in Table~\ref{points}. A pure imaginary value
$\ii y$, $y\ge1$, is in the usual modular fundamental domain and has a
unique such imaginary representative. The listed tuples are therefore
distinct under separate modular factors and torus permutations. This does
not classify all flux multiplicities at a fixed point or extra string
dualities mixing $S$ and $U$.

\begin{table}[t]
\centering
\begin{tabular}{c l c}
\toprule
$k$ & $(\Imn S;\Imn U_1,\Imn U_2,\Imn U_3)$ & $Q_{\min}$ at each point\\
\midrule
1 & $(1;1,1,1)$ & 8\\
2 & $(2;1,1,1)$, $(1;1,1,2)$ & 16\\
  & $(2;1,2,2)$, $(4;2,2,2)$ & 16\\
3 & $(3;1,1,1)$, $(1;1,1,3)$ & 24\\
  & $(3;1,3,3)$, $(9;3,3,3)$ & 24\\
\bottomrule
\end{tabular}
\caption{Nine reduced moduli point types attained within capacity 28.
The minima are over the full stated even lattice at each point, before
imposing capacity. They are not total landscape counts.}
\label{points}
\end{table}

\section{Exact pointwise minima}
Let $C_e$ be the real coefficients of the 16 multiaffine monomials of $W$,
with exponent vectors $e\in\{0,1\}^4$. The map from flux to $C$ is a signed
permutation, so $C\in(2\Z)^{16}$. Define
\begin{equation}
\mathcal A(x,y)=\frac1y\begin{pmatrix}1&x\\x&x^2+y^2\end{pmatrix},\quad
\mathcal J=\begin{pmatrix}0&-1\\1&0\end{pmatrix}.
\end{equation}
Direct tensor algebra gives the normalized positive identity
\begin{align}
E&=\tfrac14 C^T\bigotimes_A\mathcal A(\Re z_A,y_A)C,
\quad Q=\tfrac14C^T\mathcal J^{\otimes4}C,\\
E-Q&=e^K\sum_A4y_A^2|D_AW|^2.
\label{energy}
\end{align}
One derivation uses $u=(1,z)^T$ and
$uu^\dagger+\bar u u^T=2y\mathcal A$,
$uu^\dagger-\bar u u^T=2\ii y\mathcal J$. The total and alternating sums
of $|W|^2$ over the 16 patterns of conjugating individual variables produce
the two quadratic forms in \eqref{energy}. Because $C$ is real,
three-variable flips are conjugate to complementary one-variable flips;
finally $2y_AD_AW=\ii W(z_A\mapsto\bar z_A)$ establishes the last equality.

At a pure imaginary point, let $\bar e=\mathbf1-e$ and
$w_e=\prod_Ay_A^{2e_A-1}$. The contribution of each complementary pair to
$E-Q$ is
\begin{equation}
\frac14\left(\sqrt{w_e}C_e-epsilon_eC_{\bar e}/\sqrt{w_e}\right)^2,
\qquad\epsilon_e=(-1)^{\sum_Ae_A}.
\end{equation}
For $w_e=u/v>0$ in lowest terms, an F-term zero must therefore have
\begin{equation}
(C_e,C_{\bar e})=2r(v,\epsilon_eu),\quad r\in\Z,
\qquad Q_e=2uvr^2.
\label{pair}
\end{equation}
This is a saturated pairwise integral parametrization, not a flux scan.

\begin{lemma}[Necessary number of complementary pairs]
At a pure imaginary F-term zero, full real rank requires at least four
nonzero complementary monomial pairs in $W$.
\end{lemma}
\begin{proof}
If there are $r\le3$ pairs, choose real scaling weights $\ell_A$ satisfying
$(2e_j-\mathbf1)\cdot\ell=0$ for each supported pair. There are at least
$4-r$ independent choices. Every supported monomial has weight
$w=\frac12\sum_A\ell_A$, so under $z_A\mapsto e^{t\ell_A}z_A$,
\begin{equation}
W\mapsto e^{tw}W,\qquad D_AW\mapsto e^{t(w-\ell_A)}D_AW.
\end{equation}
The upper half-planes are preserved and the same fixed-flux zero has a
continuous orbit. Its independent tangent vectors $\delta z_A=\ell_Az_A$
lie in the real Jacobian kernel. Discrete modular identifications do not
remove these infinitesimal directions.
\end{proof}

\begin{proposition}
For every integer $k>1$, the four point types
\begin{equation}
(k;1,1,1),\quad(1;1,1,k),\quad(k;1,k,k),\quad(k^2;k,k,k)
\end{equation}
have exact pointwise minimum $Q_{\min}=8k$ in the stated even lattice.
At $(1;1,1,1)$ the minimum is eight.
\end{proposition}
\begin{proof}
At the unit point each primitive nonzero pair in \eqref{pair} costs two.
At any of the four other types, $w_e$ is an odd integral power of $k$;
hence $2uv\ge2k$. Four supported pairs therefore cost at least $8k$.
For example at $(k^2;k,k,k)$ their primitive costs are $2k^5$ once,
$2k^3$ three times and $2k$ four times. The divisor witnesses and their
integral modular images attain the bound with full rank, proving equality.
\end{proof}

The positivity and integral-pair analysis apply to the entire even lattice
at each specified point, not only to the positive divisor branch. They do
not give a global lower bound over all moduli. In particular capacity four
fails at these points, but a global no-vacuum theorem at capacity four is
not inferred.

\section{Scope and reproducibility}
All masses in \eqref{spectrum} acquire the common positive factor
$e^{K_K}$ when the unstabilized K\"ahler sector is restored. In the displayed
normalization $m_{3/2}^2=10k$; the mass ratios are fixed, with no claimed
tunable hierarchy. Since $W\ne0$, these are no-scale CS/axiodilaton critical
points, not fully supersymmetric four-dimensional Minkowski vacua.
For example $m_{\min}^2/m_{3/2}^2=(28-12\sqrt5)/10$ is independent of the
common volume factor.

The exact periods avoid a large-complex-structure truncation in this model.
The $g_s=1/9$ construction improves the coupling relative to the earlier
$S=2\ii$ anchor, but establishes neither corrected K\"ahler stabilization nor
a localized warped ten-dimensional solution. Open-string fields, loop and
alpha-prime corrections are not stabilized or bounded here. The reference
to inherited geometric capacity 28 is conditional; no new geometric audit,
global coupling optimum or universal tadpole statement is supplied.

The packet command \texttt{python reproduce.py} runs without network or old
project paths. It verifies the symbolic two-integer and weak-coupling
families, all 33 divisor assignments, the nine reduced point types, their
72 primitive pair costs, and a nontrivial modular image counted as the same
orbit. It recomputes the signs, charge normalization, full real Jacobians
and canonical spectra. Source hashes identify earlier completed results;
the new portable output is distinguished from those inherited records.

\appendix
\section{Seed calculation}\label{seed}
In the interleaved ordering defined above, the real seed Jacobian is
\begin{equation}
J_{F,0}=\begin{pmatrix}
-1&-3& 1&-1& 1&-1& 1&-1\\
-3& 1& 1& 1& 1& 1& 1& 1\\
 1&-1&-1&-3& 1& 1& 1& 1\\
 1& 1&-3& 1&-1& 1&-1& 1\\
 1&-1& 1& 1&-1&-3& 1& 1\\
 1& 1&-1& 1&-3& 1&-1& 1\\
 1&-1& 1& 1& 1& 1&-1&-3\\
 1& 1&-1& 1&-1& 1&-3& 1
\end{pmatrix}.
\end{equation}
It follows from
$\partial_{x_B}F_A=W_{AB}+\ii W_B/(2y_A)$ and
$\partial_{y_B}F_A=\ii\partial_{x_B}F_A
-\delta_{AB}\ii W/(2y_A^2)$. At $y_A=1$, \eqref{hessian} gives
$H_0=J_{F,0}^TJ_{F,0}$. The script records both matrices and the factored
characteristic polynomial.

\section*{Chronology and attribution disclosure}
Earlier work was Fable 5-led; July 2026 onward results were predominantly
produced by GPT models, as reported by the operator. DeepSeek and multiple
GPT models also contributed cross-review and manuscript preparation. This
is reported chronology, not a per-result authorship allocation. This note
and its closure calculation are GPT-produced. Human author and affiliation
fields are blank, and human expert review has not occurred. Historical
source dates, the present delivery date of 29 October 2026, and demonstrable
public disclosure are distinct; a private commit does not establish public
priority.
\begin{quote}
``I genuinely do not know how authorship, intellectual credit, or ownership should be attributed for these results.''
\end{quote}

\bibliographystyle{unsrtnat}
\bibliography{references}
\end{document}
